%! Mean, Variance, and Covariance — Unit 8, Lesson 8.4 — Exit Ticket (Answer Key)
\documentclass[10pt]{article}
\usepackage{linalg-article}
\usepackage{linalg-key}   % pulls in -boxes; do NOT also load -boxes
\newcommand{\TallMath}[1]{$\displaystyle #1\rule[-1.4em]{0pt}{3.2em}$}

\begin{document}

\pageheader{Unit 8, Lesson 8.4}{Exit Ticket --- Answer Key}
\namedateperiod

\vspace{0.15in}

No notes. Show your reasoning.

\begin{enumerate}[leftmargin=1.8em, itemsep=16pt]
  \item \textbf{Center and spread.} For the data $3,5,7,9$, find the mean $\mu$ and the variance
        $\sigma^2=\frac1N\sum(x_i-\mu)^2$.
        \ansline{$\mu=\frac{3+5+7+9}{4}=6$. Deviations $-3,-1,1,3$, so $\sigma^2=\frac{9+1+1+9}{4}=5$ (and $\sigma=\sqrt5\approx2.2$).}
  \item \textbf{Read a covariance matrix.} Two features have covariance matrix
        $V=\left[\begin{smallmatrix}5&3\\3&5\end{smallmatrix}\right]$.
        (a) What is the variance of each feature? \; (b) The off-diagonal entry is $+3$: in one sentence, how do
        the two features move together? \; (c) What is the total variance (the trace)?
        \ansline{(a) Each feature has variance $5$ (the diagonal entries). (b) The covariance is \emph{positive}, so the features rise and fall \emph{together} --- above-average in one tends to go with above-average in the other. (c) Total variance $=$ trace $=5+5=10$.}
  \item \textbf{Explain.} In one sentence, why is a covariance matrix always \emph{symmetric}?
        \ansline{Because the covariance of $x$ with $y$ equals the covariance of $y$ with $x$ --- the product $(x-\mu_x)(y-\mu_y)$ is the same either way --- so the off-diagonal entries match and $V=V^{\top}$.}
\end{enumerate}

\begin{teachernote}
Sort responses into ``got it / arithmetic slip / concept slip.'' Item~1 checks the mean ($6$) and variance
($5$) mechanics --- watch for squaring the sum instead of summing the squares. Item~2 checks reading a covariance
matrix: variances on the diagonal, the \emph{sign} of the covariance (positive $=$ move together), and the trace
as total variance. Item~3 is the target idea: symmetry from $\sigma_{xy}=\sigma_{yx}$. Since this is the last
lesson, close by naming the arc --- data becomes a symmetric matrix whose eigenvectors (Unit~7 PCA) are the
directions of greatest spread.
\end{teachernote}

\end{document}
